\section{Задание}

\noindent \textbf{Создать базу знаний «Предки»}, позволяющую наиболее эффективным способом (за меньшее количество шагов, что обеспечивается меньшим количеством предложений БЗ - правил), и используя разные варианты (примеры) простого вопроса, (указать: какой вопрос для какого варианта) определить:

\begin{enumerate}[label=\arabic*., itemsep=0pt, parsep=0pt]
	\item по имени субъекта определить всех его бабушек (предки 2-го колена),
	\item по имени субъекта определить всех его дедушек (предки 2-го колена),
	\item по имени субъекта определить всех его бабушек и дедушек (предки 2-го колена),
	\item по имени субъекта определить его бабушку по материнской линии (предки 2-го колена),
	\item по имени субъекта определить его бабушку и дедушку по материнской линии (предки 2-го колена).
\end{enumerate}

\noindent Минимизировать количество правил и количество вариантов вопросов. Использовать конъюнктивные правила и простой вопрос. Для одного из вариантов ВОПРОСА задания 1 составить таблицу, отражающую конкретный порядок работы системы.\\

\noindent \textbf{2. Дополнить базу знаний правилами, позволяющими найти}

\begin{enumerate}[label=\arabic*., itemsep=0pt, parsep=0pt]
	\item \textbf{Максимум из двух чисел}
	      \begin{enumerate}[label=\alph*), ref=\alph*, itemsep=0pt, parsep=0pt]
		      \item[а)] без использования отсечения,
		      \item[в)] с использованием отсечения;
	      \end{enumerate}
	\item \textbf{Максимум из трех чисел}
	      \begin{enumerate}[label=\alph*), ref=\alph*, itemsep=0pt, parsep=0pt]
		      \item[а)] без использования отсечения,
		      \item[в)] с использованием отсечения;
	      \end{enumerate}
\end{enumerate}

\noindent Убедиться в правильности результатов.\\
Для каждого случая пункта 2 обосновать необходимость всех условий тела.\\
Для одного из вариантов ВОПРОСА и каждого варианта задания 2 составить таблицу, отражающую конкретный порядок работы системы.

\vspace{0.5cm}
\noindent Форма таблицы:
\vspace{0.2cm}

\noindent
\renewcommand{\arraystretch}{1.5}
\begin{tabularx}{\textwidth}{|p{1.2cm}|X|X|X|}
	\hline
	№ шага & Состояние резольвенты, и вывод: дальнейшие действия (почему?) & Для каких термов запускается алгоритм унификации: Т1=Т2 и каков \textbf{результат} (и подстановка) & Дальнейшие действия: прямой ход или откат (почему и к чему приводит?) \\
	\hline
	1....  & ...                                                           &                                                                                                    & Комментарий, вывод...                                                 \\
	\hline
	2      & ...                                                           &                                                                                                    & ...                                                                   \\
	\hline
\end{tabularx}

\vspace{0.3cm}
\noindent Т.к. резольвента хранится в виде стека, то состояние резольвенты требуется отображать в столбик: вершина -- сверху! Новый шаг надо начинать с нового состояния резольвенты!

\vspace{0.6cm}
\noindent Для одного из вариантов ВОПРОСА и конкретной БЗ составить таблицу, отражающую конкретный порядок работы системы, с объяснениями: очередная проблема на каждом шаге и метод ее решения; каково новое текущее состояние резольвенты, как получено; какие дальнейшие действия? (Запускается ли алгоритм унификации? Каких термов? Почему этих?) ; вывод по результатам очередного шага и дальнейшие действия.
